% Phiếu theo dõi tiến độ thực hiện, biên dịch: xelatex PhieuTheoDoiTienDo.tex (2 lần)
\documentclass[12pt]{extarticle}
\usepackage[a4paper,left=2.5cm,right=2cm,top=2cm,bottom=2cm]{geometry}
\usepackage{fontspec}
\setmainfont{Times New Roman}
\usepackage{longtable,array}
\usepackage[table]{xcolor}
\usepackage{fancyhdr}
\pagestyle{fancy}\fancyhf{}\renewcommand{\headrulewidth}{0pt}
\fancyfoot[R]{\thepage}
\setlength{\parindent}{0pt}
\renewcommand{\arraystretch}{1.15}
\newcolumntype{C}[1]{>{\centering\arraybackslash}p{#1}}
\newcolumntype{L}[1]{>{\raggedright\arraybackslash\small}p{#1}}
\newcommand{\YL}{\textbf{Yến Linh:} }
\newcommand{\ML}{\textbf{Minh Lý:} }
\newcommand{\CN}{\textbf{Cả nhóm:} }

\begin{document}

TRƯỜNG ĐẠI HỌC TÔN ĐỨC THẮNG\\
\hspace*{1.2cm}\textbf{KHOA CÔNG NGHỆ THÔNG TIN}

\vspace{0.8cm}
\begin{center}
{\Large\bfseries PHIẾU THEO DÕI TIẾN ĐỘ THỰC HIỆN\\KLTN/ĐATN}
\end{center}
\vspace{0.4cm}

\begingroup\setlength{\parskip}{6pt}
Họ tên sinh viên 1 (nhóm trưởng): Sử Thị Yến Linh \hfill MSSV: 52300218

Họ tên sinh viên 2: Lê Minh Lý \hspace{1cm} MSSV: 52300220

Ngành học: Mạng máy tính và truyền thông dữ liệu

Địa chỉ liên lạc:\dotfill Đơn vị:\makebox[4cm]{\dotfill}

Họ tên GVHD: TS. Hồ Thị Linh

Mã đề tài: ĐA069

Tên đề tài KLTN: Xây dựng mô hình khuyến nghị lai kết hợp nhân tử hoá ma trận và
mạng lưới thần kinh sâu
\endgroup

\vspace{0.4cm}
\begin{longtable}{|C{2.4cm}|L{5.1cm}|L{4.6cm}|C{2.1cm}|}
\hline
 & \multicolumn{2}{c|}{\textbf{Khối lượng}} & \\ \cline{2-3}
\textbf{Tuần/Ngày} & \multicolumn{1}{c|}{\textbf{Đã thực hiện}} & \multicolumn{1}{c|}{\textbf{Tiếp tục thực hiện}} & \textbf{GVHD ký} \\ \hline
\endfirsthead
\hline
\textbf{Tuần/Ngày} & \multicolumn{1}{c|}{\textbf{Đã thực hiện}} & \multicolumn{1}{c|}{\textbf{Tiếp tục thực hiện}} & \textbf{GVHD ký} \\ \hline
\endhead

Tuần 1–2\newline 24/08 –\newline 06/09/2026 &
\YL Nghiên cứu MF, BPR và kiến trúc NCF/NeuMF; xác định bài toán khuyến nghị Top-K từ phản hồi ẩn; viết đề cương sơ bộ; xây dựng quy trình dữ liệu bản đầu (gộp cặp khách–sản phẩm, lọc k-core, chia tập huấn luyện/kiểm thử).\newline
\ML Khảo sát bộ dữ liệu và độ đo (HR, NDCG, Recall, Precision); viết đề cương chi tiết; cài đặt GMF, MLP, NeuMF và các mô hình đối chứng cổ điển; phân tích kết quả bản đầu. &
\YL Thiết kế lại hệ thống (bản 2); xây dựng bộ đọc dữ liệu H\&M.\newline
\ML Xây dựng module đánh giá xếp hạng đầy đủ. & \\ \hline

Tuần 3–4\newline 07/09 –\newline 20/09/2026 &
\YL Thiết kế bản 2 cấu hình hoá; bộ đọc dữ liệu H\&M (31,8 triệu giao dịch), lấy mẫu theo khách hàng để giữ trọn lịch sử mua; viết Chương 1–2, dựng khung báo cáo LaTeX.\newline
\ML Xây dựng module đánh giá xếp hạng toàn bộ và xếp hạng trên tập lấy mẫu, đối chiếu hai cách; bổ sung mô hình hợp nhất sớm; tổng hợp kết quả nhiều seed. &
\YL Chia dữ liệu theo mốc thời gian; xây dựng đặc trưng sản phẩm, khách hàng.\newline
\ML Xây dựng quy trình đánh giá có đăng ký trước, tinh chỉnh trên tập xác thực. & \\ \hline

Tuần 5–6\newline 21/09 –\newline 04/10/2026 &
\YL Chuyển dữ liệu chính sang H\&M, chia theo mốc thời gian chung, lọc k-core chỉ trên dữ liệu trước mốc kiểm thử; xây dựng đặc trưng sản phẩm, khách hàng, thời gian; thiết kế mô hình đề xuất NeuMF-F và LateFusion-F; viết Chương 3.\newline
\ML Đăng ký trước kế hoạch thực nghiệm; tinh chỉnh chỉ trên tập xác thực, khoá tập kiểm thử; huấn luyện lại, đánh giá cuối; viết Chương 4 bản đầu.\newline
\CN Báo cáo tiến độ với GVHD ngày 04/10. &
\CN Sửa theo góp ý của GVHD.\newline
\YL Thiết kế lại dữ liệu; chọn bộ dữ liệu kiểm chứng thứ hai.\newline
\ML Đánh giá cuối 5 seed trên H\&M. & \\ \hline

Tuần 7–8\newline 05/10 –\newline 18/10/2026 &
\YL Gộp sản phẩm theo mẫu thiết kế, tạo mẫu kiểm định mới; chọn Amazon Beauty làm dữ liệu kiểm chứng; xây dựng mô hình lai truyền thống Hybrid-CBCF; hoàn thiện Chương 2–3.\newline
\ML Đánh giá cuối 5 seed trên H\&M, kiểm chứng 5 seed trên Amazon Beauty; bảng kết quả ở nhiều giá trị K; viết phần kiểm chứng Amazon (Chương 4–5); demo chạy trên cả hai bộ dữ liệu, khách mới nhận danh sách sản phẩm bán chạy. &
\CN Báo cáo giữa kỳ với GVHD (23/10).\newline
\YL Trình bày lý thuyết, dữ liệu, mô hình đề xuất.\newline
\ML Slide phần thực nghiệm, kết quả, demo. & \\ \hline

Tuần 9–10\newline 19/10 –\newline 01/11/2026 &
\CN Chốt phạm vi, không thêm mô hình mới (20/10); báo cáo giữa kỳ với GVHD (23/10), ghi nhận góp ý.\newline
\YL Chỉnh văn phong, định dạng Chương 1–5 và tóm tắt theo góp ý; vẽ lại các sơ đồ báo cáo; thêm hình xử lý dữ liệu, hình chia theo thời gian và mục Thiết lập thực nghiệm.\newline
\ML Sửa lỗi mã nguồn và demo; vẽ lại hình kiến trúc NeuMF-F; cập nhật bảng kết quả ở K = 10, 15, 20; demo so sánh NeuMF-F với mô hình đối chứng được chọn. &
\YL Rà soát toàn văn báo cáo.\newline
\ML Kiểm lại toàn bộ số liệu; làm slide bảo vệ và poster. & \\ \hline

\multicolumn{1}{|l|}{\parbox[t]{2.4cm}{Kiểm tra\newline giữa kỳ}} &
\multicolumn{3}{l|}{\parbox[t]{11.8cm}{Đánh giá khối lượng hoàn thành\makebox[2cm]{\dotfill}\%\par
\hfill được tiếp tục/không tiếp tục thực hiện LVTN\vspace{4pt}}} \\ \hline

Tuần 11–12\newline 02/11 –\newline 15/11/2026 &
\CN Họp GVHD ngày 09/11, thống nhất màu hình, cách tính thời gian huấn luyện và slide bảo vệ.\newline
\YL Rà soát toàn văn, thống nhất thuật ngữ, trích dẫn; biên dịch bản PDF; kiểm lại số liệu và chỉnh demo (09/11); sửa báo cáo, sơ đồ theo góp ý GVHD (12/11).\newline
\ML Kiểm lại số liệu báo cáo so với kết quả đã lưu; vẽ lại hình kết quả; làm slide và poster theo mẫu Khoa; chỉnh giao diện demo (mã khách KH0001…) và poster (10/11).\newline
\textbf{Kết quả chính:} NeuMF-F đạt NDCG@10 = 0,04003 trên H\&M, cao hơn NeuMF chỉ dùng ID 43,8\% (p Holm < 0,001), tương đương nhánh GMF-F; LateFusion-F cao nhất trên cả hai bộ dữ liệu (0,04093 trên H\&M, 0,03747 trên Amazon Beauty). &
\YL Xuất bản Word của báo cáo.\newline
\CN Kiểm tra đạo văn; gửi báo cáo, slide, poster cho GVHD góp ý. & \\ \hline

Tuần 13–14\newline 16/11 –\newline 29/11/2026 &
\YL Xuất bản Word của báo cáo, cập nhật bản PDF, slide, poster (16/11).\newline
\CN Kiểm tra đạo văn bằng Turnitin: 11\% (16/11). &
\CN Gửi bản PDF, Word, slide, poster cho GVHD góp ý; sửa theo góp ý.\newline
\YL Luyện trình bày phần tổng quan, lý thuyết, dữ liệu, mô hình; cập nhật tài liệu vấn đáp.\newline
\ML Luyện trình bày phần thực nghiệm, kết quả; chạy thử demo với các khách ví dụ. & \\ \hline

Tuần 15\newline 30/11 –\newline 04/12/2026 & &
\CN Nộp báo cáo cuối (bản PDF và Word).\newline
\ML Đóng gói mã nguồn, demo và mô hình đã lưu.\newline
\YL Kiểm tra lần cuối hình thức báo cáo. & \\ \hline
\end{longtable}

\end{document}
